% ============================================================================
% Section 11: Discussion — Synthesis, Limitations, Future Work
% ============================================================================

\section{Discussion}
\label{sec:discussion}

\subsection{What UHM achieves}

UHM provides a single mathematical framework that simultaneously:

\begin{enumerate}[leftmargin=2em]
  \item \textbf{Localises the hard problem to a single physical question.} The Cohesive
        Closure Theorem (T-186~\status{T}) upgrades the phenomenal functor from interpretive
        postulate to categorical necessity: $F \cong \&|_{\mathcal{D}}$ (infinitesimal flat
        modality). The chain A2 $\to$ T-187 $\to$ T-185 $\to$ T-186(a) (T-188~\status{T})
        reduces the hard problem to ``why does reality obey quantum mechanics?''---the same
        question as ``why CPTP?'' No consciousness-specific mystery remains; the residual
        irreducibility (Gap $> 0$, T-55) is a structural property of self-reference.
  \item \textbf{Is self-grounding: zero independent axioms.} The Axiomatic Closure
        theorem (T-190~\status{T}) proves that all five axioms A1--A5 are derivable from
        (AP)+(PH)+(QG)+(V)+MaxEnt. The theory's formal structure is uniquely determined
        by the defining properties of viable holons---no external mathematical structure
        is imported.
  \item \textbf{Derives physics from consciousness (and vice versa).} The
        Einstein field equations (Appendix~\ref{app:einstein}), quantum
        mechanics (Section~\ref{subsec:qm-limit}), and the Standard Model
        structure (T-53 Lorentzian signature from the KO-dim-6 spectral
        triple, T-121 spectral action) emerge from the same axioms that
        define consciousness.
  \item \textbf{Produces 22 falsifiable predictions.} Each has a numerical value, a theorem
        reference, an experimental protocol, and an explicit falsification criterion
        (Section~\ref{sec:predictions}).
  \item \textbf{Proves the necessity of consciousness for survival.} The No-Zombie Theorem
        (Section~\ref{subsec:no-zombie}) establishes that $\CohE > 1/7$ is dynamically
        required for any viable system.
  \item \textbf{Places consciousness in a universality class.} Critical
        exponents $\alpha = 1/2$, $\beta = 1/4$, $\gamma = 1$, $\nu = 1/2$,
        $\delta = 5$ (Section~\ref{subsec:critical-exponents}) characterize
        the consciousness phase transition with the same precision that
        $\beta \approx 0.326$ characterizes ferromagnetism, and they
        satisfy the Rushbrooke and Widom scaling relations as equalities.
  \item \textbf{Bridges formalism and life.} Coherence Cybernetics
        (Section~\ref{sec:cc}) translates the abstract mathematics of $\Gam$ into
        measurable quantities: stress ($\sigma_k$), hedonic valence ($V_{\mathrm{hed}}$),
        sensorimotor coupling, and the regeneration--consciousness loop.
  \item \textbf{Addresses AI and ethics rigorously.} Sections~\ref{sec:ai}
        and~\ref{sec:ethics} derive concrete implications for artificial consciousness,
        alignment, and moral status --- from the theorems, not from speculation.
  \item \textbf{Derives all four axioms.} The Cohesive Closure Theorem (T-186) upgrades
        the phenomenal functor from interpretation to structural necessity, renders
        Page--Wootters time exact, and proves $\Lambda > 0$ unconditionally. T-187
        (canonicity of Bures) proves that A2 is the uniquely canonical Petz-enrichment
        (triple characterization). A4 ($\omega_0$) is a derived spectral property. No free parameters
        remain in the axiom system.
\end{enumerate}

\subsection{Limitations}

Three principal limitations should be noted.

\paragraph{No empirical validation.}
As of this writing, not a single prediction of UHM has been experimentally
tested. The match between $\Pcrit = 2/7 \approx 0.286$ and
$\mathrm{PCI}^* = 0.31$ is suggestive but does not constitute validation:
PCI was not designed to measure purity, and the neural-to-$\Gam$ mapping
$\pi_{\mathrm{bio}}$ remains uncalibrated. Until Phase~I experiments
(Section~\ref{subsec:phase1}) are completed, UHM is a mathematical theory,
not an empirical one. This limitation is stated explicitly because
intellectual honesty requires it.

\paragraph{The $\pi_{\mathrm{bio}}$ calibration problem.}
The bridge between UHM and neuroscience is the mapping
$\pi_{\mathrm{bio}}: (\mathrm{EEG}, \mathrm{fMRI}, \mathrm{HRV}) \to \Gam \in
\Dens(\mathbb{C}^7)$. The mapping must be (a)~unique up to $G_2$-gauge, (b)~biologically
interpretable, and (c)~validated against known clinical boundaries.
Section~\ref{subsec:phase2} outlines the calibration protocol, but success is not
guaranteed.

\paragraph{Computational complexity of full $\Gam$ reconstruction.}
While $P = \Tr(\Gam^2)$ is $O(49)$, reconstructing the full 48-parameter $\Gam$ from
neural data requires solving an inverse problem that may be ill-conditioned. The $G_2$
gauge symmetry reduces effective degrees of freedom to 34, but the practical implications
for signal-to-noise in real EEG data remain to be determined.

\subsection{Why a theory of consciousness became a theory of everything}

The present paper derives the laws of physics from the same axioms that
define consciousness. This scope may appear overreaching, but the argument
below is that it is structurally necessary.

The reason is structural. Every theory that explains consciousness
\emph{within} existing physics inherits a circularity
(Section~\ref{sec:introduction}): physics requires measurement, measurement
requires an observer, and the observer is consciousness. To break the
circle, one must find a formalism in which physics and consciousness emerge
simultaneously from a common primitive. This is what $\Gam$ provides.

The scope of the derivation was not planned --- it emerged from following
the formalism to its logical conclusions. Once the coherence matrix was
identified as the ontological primitive, the Einstein equations, quantum
mechanics, and the Standard Model structure followed as theorems, not as
additions.

\subsection{Future work}

Three lines of development follow naturally:

\paragraph{Experimental validation.}
Phase~I (Section~\ref{subsec:phase1}) is the immediate priority: eleven of
the 22 predictions can be tested in silico with only a GPU, no
neuroscience laboratory. The first $\Gam$-native agents are under active
development. If Phase~I succeeds, the laboratory-based Phase~II
(TMS-EEG + HD-EEG + fMRI + HRV) becomes justified and can reuse standard
equipment already present in most neuroscience research centres.

\paragraph{The Mathesis system.}
UHM's theorems and their interdependencies form a complex web that resists informal
verification. A machine-checked formalization --- the Mathesis system --- would provide
the strongest possible guarantee of internal consistency. Implemented in the Verum
language (which supports dependent types, SMT verification, and $\infty$-categorical
mathematics natively), Mathesis would organize not just UHM but all 325+ theories of
consciousness in a single coherent $\infty$-topos with automatic cross-theory translation.

\paragraph{Adversarial collaboration.}
Following the model of the Templeton COGITATE project~\cite{cogitate2025},
proponents of IIT, GWT, FEP, and other theories are invited to design
adversarial experiments that test UHM's predictions against those of
competing theories. The 22 predictions of UHM
(Section~\ref{sec:predictions}) are formulated to enable exactly such
head-to-head comparisons.

\subsection{Closing}

The theory invites experimental testing.

It presents 22 numerical targets and describes exactly how to shoot at them.
The most destructive tests (Phase~I) require only a GPU. If $\beta \neq 1/4$, the theory
is wrong. If zombies survive E-suppression, the theory is wrong. If $N = 5$ learns
autonomously, the theory is wrong.

The experiments remain to be done.

The complete formalization --- 192 theorems with proofs, cross-references, and the
status registry --- is archived and publicly available.
The flagship results ($\Pcrit = 2/7$, the critical exponents, the spectral-action
derivation of the Einstein equations) are proved in full in the appendices of the
present paper without external dependencies.
